% K41_Variational_Minimiser_v1_de.tex
% Paper A: Das K41-Spektrum als eindeutiger variationeller Minimierer
%          einer skalenaufgelösten freien Energie
% Author: Lukas Geiger, 2026

\documentclass[12pt,a4paper]{article}
\usepackage[utf8]{inputenc}
\usepackage[T1]{fontenc}
\usepackage[ngerman]{babel}
\usepackage{lmodern}
\usepackage{microtype}
\usepackage{amsmath,amssymb,amsthm,mathrsfs}
\usepackage{mathtools}
\usepackage{enumitem}
\usepackage{array}
\usepackage{xcolor}
\usepackage{geometry}
\geometry{a4paper, margin=2.5cm}
\setlength{\emergencystretch}{3em}
\usepackage{tikz}
\usetikzlibrary{arrows.meta,positioning,calc}
\usepackage{xurl}
\usepackage[unicode=true]{hyperref}
\hypersetup{
  colorlinks=true,
  linkcolor=blue!70!black,
  citecolor=green!50!black,
  urlcolor=blue!70!black,
  pdftitle={Das K41-Spektrum als eindeutiger variationeller Minimierer einer skalenaufgelösten freien Energie},
  pdfauthor={Lukas Geiger},
  pdfsubject={Variationelle Charakterisierung des Kolmogorov-K41-Energiespektrums},
  pdfkeywords={K41-Spektrum, Turbulenz, Variationsprinzip, freie Energie, Kolmogorov-Skalierung}
}
\providecommand*{\HyperDestNameFilter}[1]{de.#1}
\Urlmuskip=0mu plus 2mu\relax
\newcolumntype{L}[1]{>{\raggedright\arraybackslash}p{#1}}

\theoremstyle{definition}
\newtheorem{definition}{Definition}[section]
\theoremstyle{plain}
\newtheorem{theorem}[definition]{Theorem}
\newtheorem{lemma}[definition]{Lemma}
\newtheorem{proposition}[definition]{Proposition}
\newtheorem{corollary}[definition]{Korollar}
\newtheorem{conjecture}[definition]{Vermutung}
\newtheorem{axiom}[definition]{Axiom}
\theoremstyle{remark}
\newtheorem{remark}[definition]{Bemerkung}

% ---------------------------------------------------------------------------
\title{Das K41-Spektrum als eindeutiger variationeller Minimierer\\
einer skalenaufgelösten freien Energie}

\author{Lukas Geiger\\
\textit{Unabhängiger Forscher, Bernau, Deutschland}\\
ORCID: \href{https://orcid.org/0009-0005-7296-1534}{0009-0005-7296-1534}}

\date{Mai 2026}

\begin{document}
\maketitle

% ===================================================================
\begin{abstract}
Wir etablieren ein referenznormalisiertes Variationsprinzip, welches das
Kolmogorov-$k^{-5/3}$-Energiespektrum als eindeutigen Minimierer in der
Energievariable innerhalb der angegebenen K41-normalisierten zulässigen
Klasse charakterisiert. Ausgehend von einer
Littlewood--Paley-Zerlegung des Geschwindigkeitsfeldes definieren wir ein
Freie-Energie-Funktional $\mathcal{F}[\mathbf{E}]$ basierend auf der
Kullback--Leibler-Divergenz zwischen der Schalen-Energieverteilung und einem
K41-Referenzgleichgewicht. Wir führen die K41-normalisierte Klasse zulässiger
Energiefluss-Operatoren ein --- lokale, dimensionskonsistente, monotone
Schließungen (closures), die die Axiome (A1) und (A3) sowie die
Normalisierungsbedingung aus Definition~\ref{def:admissible} erfüllen ---
und beweisen, dass die K41-Verteilung der eindeutige Minimierer von
$\mathcal{F}[\mathbf{E}]$ auf der gemeinsamen zulässigen Menge von
Energieprofilen und Schließungen ist, vorausgesetzt, der Skala-zu-Skala-Energiefluss
ist konstant (Theorem~\ref{thm:K41}). Der
Minimierer ist nicht-entartet: die Hesse-Matrix von $\mathcal{F}$ bei
$\mathbf{E}^*$ ist diagonal mit Einträgen $1/E_j^* > 0$, was eine
  Schalenenergie-Steifigkeit liefert, die für dyadische Schalen wie
  $k_j^{2/3}$ wächst. Zusätzlich
  halten wir zwei Leitplanken fest: Relaxierte Flussbedingungen liefern kein
  nichttriviales Robustheitstheorem, solange das exakte K41-Profil zulässig
  bleibt, und die Minimax-Formulierung ist nur eine formale entkoppelte
  Penalty-Darstellung, kein dynamischer Zwei-Spieler-Satz.
Unter zusätzlichen Gauß-/Lognormal-Fluktuationsannahmen reproduzieren
Hesse-Fluktuationen um das Variationsminimum die Kolmogorov--Obukhov-
Exponentenformel (K62). Das Ergebnis gilt als statisches
endliches bzw. trägheitsbereichliches Variationsresultat innerhalb der
angegebenen K41-normalisierten Klasse; es ist kein dynamischer
Selektionssatz für die Navier--Stokes-Entwicklung. Anomale Dissipation
unter zusätzlichen Annahmen über die nichtlineare Kaskade wird in einer
Begleitarbeit behandelt~\cite{Geiger2026b}.
\end{abstract}

\begin{center}
\small\emph{Preprint-Hinweis. Dies ist Paper A der Turbulenz-Aufteilung. Es isoliert den statischen, referenznormalisierten K41-Variationsminimierer-Satz. Die Begleitarbeit Paper B behandelt anomale Dissipation und die Downhill-Flux-Bedingung als konditionale Kaskadenroute.}
\end{center}
\clearpage

% ===================================================================
\section{Einleitung}\label{sec:intro}

Die Turbulenz bleibt eines der herausragendsten ungelösten Probleme der klassischen
Physik und angewandten Mathematik. Trotz mehr als eines Jahrhunderts der Forschung
wurde noch keine Ableitung der statistischen Gesetze für vollständig entwickelte Turbulenz 
aus ersten Prinzipien erreicht.

\subsection{Die Kolmogorov-Theorie}

In seinen bahnbrechenden Arbeiten von 1941 postulierte Kolmogorov~\cite{K41a,K41b}, dass
im Trägheitsbereich (inertial range) von vollständig entwickelter, homogener, isotroper
Turbulenz das Energiespektrum die universelle Form
\begin{equation}\label{eq:K41}
  E(k) \;=\; C_K \,\varepsilon^{2/3}\, k^{-5/3}
\end{equation}
annimmt, wobei $\varepsilon$ die mittlere Energiedissipationsrate pro
Masseneinheit bezeichnet und $C_K \approx 1{,}5$ die Kolmogorov-Konstante ist. Diese Vorhersage,
abgeleitet aus der Dimensionsanalyse und der Hypothese der lokalen Isotropie,
wurde experimentell und numerisch über einen riesigen Bereich von
Reynolds-Zahlen hinweg bestätigt~\cite{Frisch95,Pope00}.

\subsection{Unser Beitrag}

In dieser Arbeit wählen wir einen anderen Ansatz. Anstatt das
K41-Spektrum allein aus der Dimensionsanalyse abzuleiten, zeigen wir, dass es die
\emph{im Sinn der freien Energie bevorzugte} Verteilung innerhalb einer
K41-normalisierten Konstantfluss-Vergleichsklasse ist.
Wir führen ein Freie-Energie-Funktional~$\mathcal{F}$ auf dem Raum der
skalenaufgelösten Energieverteilungen ein und beweisen:
\begin{enumerate}
  \item Das K41-Spektrum ist der \emph{eindeutige} Minimierer von $\mathcal{F}$
        unter der Bedingung eines konstanten Energieflusses, wobei die Minimierung
        gemeinsam über K41-normalisierte zulässige lokale Flussoperatoren und
        Energieprofile läuft (Theorem~\ref{thm:K41}). Das Argument setzt keine
        einzelne feste Schließung (closure) voraus, bleibt aber ausdrücklich
        referenznormalisiert.
  \item Der Minimierer ist nicht-entartet, wobei die Hesse-Matrix eine
        skalenabhängige Schalenenergie-Steifigkeit
        $1/E_j^* \propto k_j^{2/3}$ liefert, die
        kleinskalige Störungen stärker bestraft als großskalige --- ein
        statisches konvexes Signal, das mit einer abwärtsgerichteten
        Kaskaden-Tendenz vereinbar ist.
  \item Relaxierte Flussbedingungen und die Minimax-Notation werden als
        Geltungsbereich-Leitplanken festgehalten: Solange das exakte
        K41-Profil zulässig bleibt, ist das relaxierte Problem trivial, und
        die Minimax-Form ist eine entkoppelte Penalty-Darstellung statt ein
        dynamischer Zwei-Spieler-Satz
        (Propositionen~\ref{prop:dfc-slack}--\ref{prop:minimax}).
\end{enumerate}
Die dynamischen Konsequenzen der Freie-Energie-Struktur, einschließlich
anomaler Dissipation unter zusätzlichen Hypothesen zur nichtlinearen
Kaskade, werden in einem Begleitpapier entwickelt~\cite{Geiger2026b}.

\subsection{Strukturelle Klassifizierung}

Innerhalb des umfassenderen Programms der Spektraltheorie der freien
Energie~\cite{FST-Unified2026} zeigt der K41-Selektionsmechanismus die
\textbf{Pattern~A Stabilitätslogik}: (a)~Das Energiebudget in führender Ordnung ist
neutral (jedes Potenzgesetz-Spektrum ist kinematisch zulässig); (b)~Die
Bedingung konstanten Flusses erster Ordnung schränkt das Profil ein, legt es aber nicht
eindeutig fest; (c)~Die strikte Konvexität zweiter Ordnung
von~$\mathcal{F}$ wählt K41 als eindeutigen Minimierer aus. Diese dreistufige
Hierarchie wird in Abschnitt~\ref{sec:resolvent-cascade} diskutiert.

Diese Arbeit ist in sich geschlossen; alle Hauptresultate stützen sich ausschließlich auf Definitionen
und Beweise innerhalb dieses Dokuments. Verbindungen zum umfassenderen FST-Programm
werden als Kontext zitiert, sind aber logisch nicht zwingend erforderlich.

Zur Orientierung trennt Tabelle~\ref{tab:claim-status-ledger} die
satzförmige Kernaussage des Papers von diagnostischen und offenen
Brückenbausteinen. Das ist die beabsichtigte Lesereihenfolge: Der
Variationssatz ist innerhalb der angegebenen Klasse geschlossen, während
Numerik, Intermittenz und DNS nur Stützschichten oder zukünftige
Prüfgates sind.

\begin{table}[ht]
\centering
\small
\setlength{\tabcolsep}{4pt}
\begin{tabular}{|L{0.21\textwidth}|L{0.35\textwidth}|L{0.33\textwidth}|}
\hline
\textbf{Baustein} & \textbf{Evidenz- oder Darstellungsebene} & \textbf{Aktueller Status und Einschränkung} \\
\hline
Statischer Variationskern & KL-Funktional, zulässige Flussfamilie und Beweis über die gemeinsame zulässige Menge. & Bewiesen innerhalb der K41-normalisierten inertialen Schalenklasse; kein intrinsischer DNS- oder Navier--Stokes-Selektionssatz. \\
\hline
Lokale Steifigkeit & Diagonale Hesse-Matrix bei $\mathbf{E}^*$ mit Schalenenergie-Skala $1/E_j^* \propto k_j^{2/3}$. & Bewiesen als lokale Konvexität; liefert statische Steifigkeit, keine Rückkehrdynamik. \\
\hline
Leitplanken-Zweige & Propositionen zu relaxiertem Fluss und Minimax-Form. & Nur Geltungsbereichskontrollen; kein positiver Robustheitssatz und kein dynamisches Spiel. \\
\hline
Synthetische Numerik & Sechs Modellspektren und $300$ zufällige Perturbationstests aus \texttt{compute\_F\_spectrum.py}. & Illustration und Audit-Check; keine unabhängige empirische Validierung. \\
\hline
K62- und Intermittenzbrücke & Gauß-/Lognormal-Fluktuationsroute um die Hesse-Skala. & Formale Modellroute; ein Kaskadenmaß und eine Strukturfunktionen-Brücke bleiben nötig. \\
\hline
DNS- und Paper-B-Brücke & Vorgeschlagenes DNS-Protokoll und Begleitpaper zur Kaskadenroute. & Zukünftige oder konditionale Ebene; hängt an DFC- und Nichtlinearitäts-Hypothesen außerhalb dieses Papers. \\
\hline
\end{tabular}
\caption{Claim-Status-Ledger für das K41-Variationsminimierer-Paper.}
\label{tab:claim-status-ledger}
\end{table}

% ===================================================================
% Figure 1: Proof chain overview
% ===================================================================
\begin{figure}[ht]
\centering
\begin{tikzpicture}[
  box/.style={draw, rounded corners, minimum width=3.2cm, minimum height=0.9cm,
              align=center, font=\small},
  arrow/.style={-{Stealth[length=5pt]}, thick},
  note/.style={font=\scriptsize\itshape, text=gray}
]
% Row 1: Axioms
\node[box, fill=blue!10] (A1) at (0,0) {(A1) Lokalität\\(A3) Monotonie};
\node[box, fill=blue!10] (LP) at (5,0) {Littlewood--Paley\\-Zerlegung};

% Row 2: Functional
\node[box, fill=orange!15] (F) at (2.5,-1.8) {Freie Energie\\$\mathcal{F}[E] = \sum_j \mathrm{KL}(E_j \| E_j^*)$};

% Row 3: Main Theorem
\node[box, fill=green!15] (T1) at (2.5,-3.6) {\textbf{Thm\,1:} K41\\eindeutiger Minimierer};

% Row 4: Extensions
\node[box, fill=yellow!15] (R) at (-0.5,-5.4) {Relaxierter Fluss\\(Prop.~\ref{prop:dfc-slack})};
\node[box, fill=yellow!15] (M) at (5.5,-5.4) {Minimax\\(Prop.~\ref{prop:minimax})};

% Arrows
\draw[arrow] (A1) -- (F) node[midway, left, note] {Def.\,\ref{def:F}};
\draw[arrow] (LP) -- (F) node[midway, right, note] {Schalenenergie};
\draw[arrow] (F) -- (T1) node[midway, right, note] {gemeinsame Min.};
\draw[arrow] (T1) -- (R) node[midway, left, note] {Geltungsbereich};
\draw[arrow] (T1) -- (M) node[midway, right, note] {Penalty};
\end{tikzpicture}
\caption{Logische Struktur des Beweises. Das Freie-Energie-Funktional
$\mathcal{F}[\mathbf{E}]$ wird aus der Littlewood--Paley-Zerlegung und 
der zulässigen Flussfamilie (Axiome A1, A3) konstruiert.
Theorem~\ref{thm:K41} (K41-Eindeutigkeit) folgt aus der gemeinsamen Minimierung.
Die relaxierte Flussbedingung und die Minimax-Verzweigung halten
Geltungsbereich-Leitplanken und eine formale Penalty-Darstellung fest, nicht
zusätzliche dynamische Theoreme.}
\label{fig:proof-chain}
\end{figure}

% ===================================================================
\section{Das skalenaufgelöste Freie-Energie-Funktional}\label{sec:functional}

\subsection{Littlewood--Paley-Zerlegung}

Sei $u \in L^2(\mathbb{T}^3)$ ein divergenzfreies Geschwindigkeitsfeld auf dem
3-Torus. Wir verwenden die Standard-Littlewood--Paley-Zerlegung
\begin{equation}\label{eq:LP}
  u \;=\; \sum_{j \ge -1} \Delta_j u,
\end{equation}
wobei $\Delta_j$ der Frequenz-Lokalisierungsoperator an der dyadischen Schale
$|\xi| \sim 2^j$ ist. Wir setzen $k_j = 2^j$ und definieren die \emph{Schalenenergie} (shell energy)
\begin{equation}\label{eq:Ej}
  E_j \;:=\; \|\Delta_j u\|_{L^2}^2.
\end{equation}
Im Trägheitsbereich übersetzt sich die K41-Vorhersage~\eqref{eq:K41} in
\begin{equation}\label{eq:Ej_star}
  E_j^* \;=\; C_K\, \varepsilon^{2/3}\, k_j^{-5/3} \,\Delta k_j,
  \qquad \Delta k_j = k_j \ln 2 \quad (\text{dyadische Schalenbreite}).
\end{equation}

\subsection{Energiefluss-Bedingung}

Der Skala-zu-Skala-Energiefluss durch die Wellenzahl $k_j$ wird über den
trilinearen Term in den Navier--Stokes-Gleichungen definiert:
\begin{equation}\label{eq:flux}
  \Pi_j \;=\; -\sum_{\ell \le j}
    \langle \Delta_\ell(u \cdot \nabla u),\; \Delta_\ell u \rangle.
\end{equation}
Im Trägheitsbereich erfüllt eine statistisch stationäre Kaskade die Bedingung
\begin{equation}\label{eq:const_flux}
  \Pi_j \;=\; \varepsilon \quad \text{für alle } j_{\mathrm{in}} \le j \le j_{\mathrm{dis}}.
\end{equation}

\subsection{Das Funktional}

\begin{definition}[Skalenaufgelöste freie Energie]\label{def:F}
Sei $\mathbf{E} = (E_j)_{j \in \mathcal{J}}$ eine Verteilung von Schalenenergien
über den Trägheitsbereich $\mathcal{J} = \{j_{\mathrm{in}}, \dots,
j_{\mathrm{dis}}\}$, mit $E_j > 0$. Sei $\mathbf{E}^* = (E_j^*)$ die
Kolmogorov-Referenzverteilung~\eqref{eq:Ej_star}. Das
\emph{skalenaufgelöste Freie-Energie-Funktional} ist
\begin{equation}\label{eq:F}
  \mathcal{F}[\mathbf{E}]
  \;:=\;
  \sum_{j \in \mathcal{J}}
  \Bigl[
    E_j \ln\!\frac{E_j}{E_j^*} \;-\; E_j \;+\; E_j^*
  \Bigr].
\end{equation}
\end{definition}

\begin{remark}
Der Summand $f(E_j) = E_j \ln(E_j/E_j^*) - E_j + E_j^*$ ist die
Kullback--Leibler-Divergenz (relative Entropie) des Paares $(E_j, E_j^*)$,
wenn sie als unnormalisierte Maße auf einem Einzelpunktraum betrachtet werden. Er erfüllt
$f(E_j) \ge 0$ mit Gleichheit genau dann, wenn $E_j = E_j^*$. Somit
gilt $\mathcal{F}[\mathbf{E}] \ge 0$ mit Gleichheit genau dann, wenn $\mathbf{E}
= \mathbf{E}^*$.
\end{remark}

\begin{figure}[ht]
\centering
\begin{tikzpicture}[scale=0.9]
% Axes
\draw[-{Stealth}] (-0.3,0) -- (7,0) node[right] {$E_j / E_j^*$};
\draw[-{Stealth}] (0,-0.3) -- (0,4.5) node[above] {$f(E_j)$};
% KL curve: f(x) = x ln(x) - x + 1, plotted for x = E_j/E_j*
\draw[thick, blue!70!black, domain=0.08:4.5, smooth, samples=80]
  plot ({\x}, {\x*ln(\x)/ln(2.718) - \x + 1});
% Minimum point
\fill[red] (1,0) circle (2.5pt);
\node[below, font=\small] at (1,-0.15) {$E_j^*$ (K41)};
% Annotation
\draw[dashed, gray] (1,0) -- (1,3.5);
\node[font=\small, text=blue!70!black, right] at (4.5,2.8) {$f = E_j\ln\frac{E_j}{E_j^*} - E_j + E_j^*$};
% Flux constraint region
\draw[thick, red!60!black, dashed] (0.3,0.2) -- (0.3,3.8);
\draw[thick, red!60!black, dashed] (3.5,0.2) -- (3.5,3.8);
\node[font=\scriptsize, text=red!60!black] at (1.9,4.1) {zulässiger Bereich};
\draw[{Stealth}-{Stealth}, red!60!black] (0.3,3.9) -- (3.5,3.9);
\end{tikzpicture}
\caption{Die Freie-Energie-Dichte $f(E_j) = E_j \ln(E_j / E_j^*) - E_j + E_j^*$
als Funktion des Schalenenergie-Verhältnisses $E_j / E_j^*$. Das eindeutige globale Minimum
bei $E_j = E_j^*$ (K41-Spektrum) ist markiert. Die Flussbedingung schränkt den
zulässigen Bereich ein, aber das Minimum liegt in der K41-normalisierten
Vergleichsklasse im Inneren und illustriert K41 als gemeinsamen
Energieminimierer.}
\label{fig:FE-landscape}
\end{figure}

\subsection{Dimensionskonsistenz}

Wir verifizieren, dass $\mathcal{F}$ dimensionskonsistent ist. Jedes $E_j$ hat die
Dimension $[L^2\,T^{-2}]$ (Energie pro Masseneinheit, integriert über ein
Spektralband). Da $E_j^*$ dieselbe Dimension hat, ist das Verhältnis
$E_j/E_j^*$ dimensionslos, der Logarithmus wohldefiniert und jeder
Summand in~\eqref{eq:F} hat die Dimension $[L^2\,T^{-2}]$.
% ===================================================================
\section{Hauptergebnisse}\label{sec:results}

\subsection{Eindeutigkeit der Kolmogorov-Skalierung}

\begin{definition}[Zulässige Flussfamilie]\label{def:admissible}
Eine Familie von Flussoperatoren $\bigl(\Pi_j[\mathbf{E}]\bigr)_{j \in
\mathcal{J}}$ heißt \emph{zulässig}, wenn jedes $\Pi_j$ Folgendes erfüllt:
\begin{enumerate}
  \item[\textup{(A1)}] \textbf{Lokalität.}\;
    $\Pi_j$ hängt nur von $E_{j-1},\, E_j,\, E_{j+1}$ ab.
  \item[\textup{(A2)}] \textbf{Dimensionskonsistenz.}\;
    $[\Pi_j] = [E_j / \tau_j]$ mit der Eddy-Turnover-Zeit (Wirbel-Umschlagszeit)
    $\tau_j = (k_j^2\,E_j)^{-1/2}$, sodass $\Pi_j$ die Form
    \begin{equation}\label{eq:closure_general}
      \Pi_j[\mathbf{E}]
      \;=\;
      k_j\,E_j^{3/2}\;\Phi_j\!\Bigl(
        \frac{E_{j-1}}{E_j},\;
        \frac{E_{j+1}}{E_j}
      \Bigr)
    \end{equation}
    besitzt, wobei $\Phi_j : (0,\infty)^2 \to (0,\infty)$ eine glatte,
    dimensionslose Funktion ist, die $\Phi_j(r^*_{-},\, r^*_{+}) = \alpha$
    für eine universelle Konstante $\alpha > 0$ erfüllt, mit
    $r^*_{\pm} = E_{j \pm 1}^* / E_j^*$.
  \item[\textup{(A3)}] \textbf{Monotonie.}\;
    $\partial \Pi_j / \partial E_j > 0$ für alle $\mathbf{E} > 0$.
\end{enumerate}
Wir bezeichnen die Menge aller zulässigen Familien mit $\mathscr{A}$.
\end{definition}

\begin{remark}[Referenznormalisierte zulässige Klasse]\label{rem:reference-normalised-class}
Die Bedingung $\Phi_j(r^*_{-},r^*_{+})=\alpha$ normalisiert die zulässige
Klasse bereits am K41-Referenzprofil. Der folgende Satz ist daher ein
relatives Selektionsresultat innerhalb dieser angegebenen K41-normalisierten
Klasse. Er behauptet nicht, dass jede feste lokale monotone Schließung K41
auswählt, und er schließt keine alternativen Schließungsfamilien aus, die um
ein anderes Referenzprofil normalisiert sind. Solche Alternativen würden
zusätzliche Skaleninvarianz-, Universalitäts- oder empirische
Zulässigkeitsannahmen benötigen.
\end{remark}

\begin{proposition}[Dimensionale Ableitung des Fluss-Vorfaktors]\label{prop:A2-derived}
Innerhalb des reduzierten skalenlokalen inertialen Schalenansatzes, in dem
bei Schale~$j$ nur $k_j$ und $E_j$ als dimensionale Größen zugelassen sind,
folgt der Vorfaktor in Axiom~\textup{(A2)} aus Dimensionsanalyse:
$\Pi_j = k_j\,E_j^{3/2}\,\Phi_j(\cdot)$.  (A2) ist damit der kanonische
Vorfaktor innerhalb dieses festgelegten Ansatzes, aber kein Satz, der
reichere Schließungsklassen mit zusätzlichen dimensionalen oder
dimensionslosen Parametern ausschließt.
\end{proposition}

\begin{proof}
Innerhalb dieses reduzierten Ansatzes erzwingt das Buckinghamsche
$\Pi$-Theorem, dass der Fluss $\Pi_j$ ein Produkt aus Potenzen der
zugelassenen dimensionalen Größen $k_j$ und~$E_j$ sein muss, multipliziert
mit einer dimensionslosen Funktion dimensionsloser Verhältnisse. Schreibe
$\Pi_j = k_j^{a}\,E_j^{b}\,\Phi_j(\ldots)$ mit
$[\Pi_j] = [L^2\,T^{-3}]$ (Energiefluss pro Masseneinheit),
$[k_j] = [L^{-1}]$ und $[E_j] = [L^2\,T^{-2}]$.

Dimensionsabgleich: $[k_j^a\,E_j^b] = L^{-a}\,L^{2b}\,T^{-2b} = L^{2b-a}\,T^{-2b}$.

Gleichsetzen mit $L^2\,T^{-3}$:
\[
2b - a = 2, \qquad 2b = 3 \quad\Longrightarrow\quad b = \tfrac{3}{2},\;\; a = 1.
\]
Daraus folgt $\Pi_j = k_j\,E_j^{3/2}\,\Phi_j$, und nach (A1) kann die
dimensionslose Funktion $\Phi_j$ nur von den lokalen Verhältnissen
$E_{j-1}/E_j$ und $E_{j+1}/E_j$ abhängen. Die Eddy-Turnover-Zeit
$\tau_j = (k_j^2\,E_j)^{-1/2}$ ist gleichermaßen die kanonische lokale
Zeitskala innerhalb derselben eingeschränkten Menge dimensionaler
Variablen. Schließungen, die Viskosität, Antriebsskalen,
Intermittenzvariablen oder weitere dimensionslose Schalenparameter
beibehalten, liegen außerhalb dieser speziellen Reduktion.
\end{proof}

\begin{remark}
Die klassische Kolmogorov--Heisenberg-Schließung $\Pi_j = \alpha\,k_j\,E_j^{3/2}$
entspricht dem Spezialfall $\Phi_j \equiv \alpha$~\cite{MY75}. Definition~\ref{def:admissible}
soll lokale Schließungsmotive wie EDQNM-artige Schließungen,
Leith-artige Diffusionsmodelle und trunkierte Triaden-Interaktionen abdecken,
wenn sie in den Schalenvariablen von Definition~\ref{def:admissible} ausgedrückt werden.
Die Familie $\mathscr{A}$ ist somit eine strukturierte lokale Vergleichsklasse
und keine einzelne Schließung.
Proposition~\ref{prop:A2-derived} trennt den reduzierten inertialen
Dimensionsansatz von den wirklich variablen Annahmen Lokalität und
Monotonie. Sie präzisiert damit den Modellierungsgehalt von (A2), statt
allen unabhängigen Modellierungsgehalt aus dem Flussvorschlag zu entfernen.
\end{remark}

\begin{theorem}[Relative Eindeutigkeit in der K41-normalisierten zulässigen Klasse]\label{thm:K41}
Sei $\varepsilon > 0$ gegeben und sei $\mathscr{A}$ die Klasse der
zulässigen Flussfamilien (Definition~\ref{def:admissible}), einschließlich
ihrer K41-Normalisierungsbedingung.
Betrachte das \emph{gemeinsame} Minimierungsproblem
\begin{equation}\label{eq:min}
  \min_{\substack{\mathbf{E} > 0,\\ \Pi \in \mathscr{A}}}
  \;\mathcal{F}[\mathbf{E}]
  \qquad \text{unter der Bedingung} \quad
  \Pi_j[\mathbf{E}] = \varepsilon
  \;\;\text{für alle } j \in \mathcal{J}.
\end{equation}
Dann gilt:
\begin{enumerate}
  \item[\textup{(i)}]
    Auf der gemeinsamen zulässigen Menge
    $\{(\mathbf{E},\Pi):\Pi\in\mathscr{A},\;
    \Pi_j[\mathbf{E}]=\varepsilon\;\text{für alle }j\}$ ist die
    K41-Verteilung
    $E_j^* = C_K\,\varepsilon^{2/3}\,k_j^{-5/3}\,\Delta k_j$
    der eindeutige Minimierer von~$\mathcal{F}$ in der Energievariable,
    und das Minimum wird
    durch das Kolmogorov--Heisenberg-Mitglied $\Phi_j \equiv \alpha$ mit
    $C_K = \bigl(\alpha\,(\ln 2)^{3/2}\bigr)^{-2/3}$ realisiert.
  \item[\textup{(ii)}]
    Der Energieminimierer ist nicht-entartet: die Hesse-Matrix
    $D^2\mathcal{F}|_{\mathbf{E}^*}$ eingeschränkt auf den Tangentialraum der
    gemeinsamen zulässigen Menge in Energierichtungen ist strikt positiv
    definit. (Dies ist eine direkte Konsequenz der strikten Konvexität der
    KL-Divergenz; ihr inhaltlicher Gehalt ist die skalenabhängige Steifigkeit
    $1/E_j^* \propto k_j^{2/3}$ für dyadische Schalenenergien,
    siehe Schritt~5b unten.)
\end{enumerate}
\end{theorem}

\begin{proof}
\textbf{Schritt 1: Gemeinsame zulässige Profile.}\;
Wir betrachten Paare $(\mathbf{E},\Pi)$ in der gemeinsamen zulässigen Menge,
für die $\Pi\in\mathscr{A}$ und $\Pi_j[\mathbf{E}]=\varepsilon$ für alle
$j$ gilt. Lokale Monotonie in $E_j$ ist für sich genommen kein globaler
Existenzsatz für das gekoppelte Schalensystem; die Existenz für eine feste
Schließung ist daher nur eine Zulässigkeitsfrage. Der Satz behauptet nicht,
dass jede feste Schließung K41 auswählt.

\textbf{Schritt 2: Untere Schranke via Konvexität.}\;
Da jeder Summand $f(E_j) = E_j\ln(E_j/E_j^*) - E_j + E_j^*$ in der Terminologie der konvexen Analysis und der
Störungsanalyse~\cite{Rockafellar70,BonnansShapiro00} strikt konvex mit globalem Minimum $f(E_j^*) = 0$ ist, haben wir
\begin{equation}\label{eq:F_lower}
  \mathcal{F}[\mathbf{E}]
  \;\ge\; 0
  \;=\; \mathcal{F}[\mathbf{E}^*],
\end{equation}
mit Gleichheit genau dann, wenn $E_j = E_j^*$ für alle~$j$.

\textbf{Schritt 3: Erreichbarkeit.}\;
Wir verifizieren, dass $\mathbf{E}^*$ tatsächlich ein stationäres Profil für
ein $\Pi \in \mathscr{A}$ ist. Wähle die Kolmogorov--Heisenberg-Schließung
$\Phi_j \equiv \alpha$. Dann ist
\[
  \Pi_j[\mathbf{E}^*]
  \;=\; \alpha\,k_j\,(E_j^*)^{3/2}
  \;=\; \alpha\,k_j\,
    \bigl(C_K\,\varepsilon^{2/3}\,k_j^{-5/3}\,k_j\ln 2\bigr)^{3/2}
  \;=\; \alpha\,C_K^{3/2}\,(\ln 2)^{3/2}\,\varepsilon.
\]
Setzen wir $C_K = \bigl(\alpha\,(\ln 2)^{3/2}\bigr)^{-2/3}$, so ergibt sich
$\Pi_j[\mathbf{E}^*] = \varepsilon$. (Der geometrische Faktor $(\ln 2)^{3/2}$,
der sich aus der dyadischen Schalenbreite ergibt, wird in die Kolmogorov-Konstante
$C_K$ absorbiert, die dadurch durch den Schließungsparameter~$\alpha$ bestimmt wird.)
Damit ist die untere Schranke~\eqref{eq:F_lower} erreicht.

\textbf{Schritt 4: Eindeutigkeit.}\;
Angenommen, $\Pi' \in \mathscr{A}$ und $\mathbf{E}' \ne \mathbf{E}^*$ bilden
ein gemeinsam zulässiges Paar, also $\Pi'_j[\mathbf{E}'] = \varepsilon$.
Dann liefert die strikte Konvexität
$\mathcal{F}[\mathbf{E}'] > 0 = \mathcal{F}[\mathbf{E}^*]$, sodass $\mathbf{E}'$
kein Minimierer von~\eqref{eq:min} sein kann. Daher ist $\mathbf{E}^*$ die
\emph{eindeutige Energielösung} des gemeinsamen Problems.

\smallskip\noindent
\textbf{Anmerkung zur logischen Struktur.}\;
Die Eindeutigkeit in Schritt~4 folgt aus der strikten Positivität der KL-Divergenz abseits der Referenz, welche für \emph{jede} Wahl
von~$\mathbf{E}^*$ gilt. Der physikalisch nichttriviale Inhalt liegt in
Schritt~3 (Erreichbarkeit): Innerhalb der K41-normalisierten Klasse gibt es
eine zulässige Schließung, die $\Pi_j[\mathbf{E}^*] = \varepsilon$
realisiert. Referenzen wie etwa $k^{-3}$ werden nicht durch KL-Konvexität
allein ausgeschlossen; sie würden eine andere Normalisierung oder stärkere
Skaleninvarianz- und Uniformitätsannahmen an die Schließungsfamilie
erfordern (siehe Bemerkungen~\ref{rem:reference-normalised-class}
und~\ref{rem:circularity}).

\textbf{Schritt 5: Nicht-Entartung.}\;
Die Hesse-Matrix $\partial^2 \mathcal{F}/\partial E_j\,\partial E_\ell
= \delta_{j\ell}/E_j$ bei $\mathbf{E}^*$ ist diagonal mit Einträgen
$1/E_j^* > 0$. Für eine \emph{feste} Schließung $\Pi \in \mathscr{A}$ stellt die
Bedingung $\Pi_j[\mathbf{E}] = \varepsilon$ genau $|\mathcal{J}|$
Gleichungen für $|\mathcal{J}|$ Unbekannte dar, was generisch isolierte
Lösungen liefert, sodass der Tangentialraum der Nebenbedingungs-Mannigfaltigkeit bei
$\mathbf{E}^{(\Pi)}$ trivial ist und Nicht-Entartung vakuös erfüllt ist.
Der wesentliche Gehalt von~(ii) betrifft das \emph{gemeinsame} Problem: Im
erweiterten Raum $(\mathbf{E}, \Pi)$ mit $\Pi$ variierend über die
(unendlichdimensionale) Familie~$\mathscr{A}$, besitzt die Einschränkungsmenge
$\{(\mathbf{E}, \Pi) : \Pi_j[\mathbf{E}] = \varepsilon\}$ tangentiale
Richtungen, entlang derer $\Pi$ variiert. Die uneingeschränkte Hesse-Matrix
$D^2_{\mathbf{E}}\mathcal{F}|_{\mathbf{E}^*}$ ist strikt positiv
definit (diagonale Einträge $1/E_j^* > 0$), daher bleibt ihre Einschränkung auf
\emph{jeden} Unterraum positiv definit.

\medskip
\noindent\textbf{Schritt 5b: Interpretation zweiter Ordnung.}\;
Da auf dyadischen Schalen
$E_j^* = C_K\varepsilon^{2/3}k_j^{-5/3}\Delta k_j
      = C_K\varepsilon^{2/3}(\ln 2)k_j^{-2/3}$ gilt,
wachsen die diagonalen Einträge $1/E_j^* \propto k_j^{2/3}$ der
Hesse-Matrix mit der Wellenzahl. Damit werden kleinskalige Störungen des
K41-Schalenenergieprofils strenger bestraft als großskalige. Diese skalenabhängige
Steifigkeit ist ein statisches konvexes Signal, das mit einer
abwärtsgerichteten Kaskaden-Tendenz vereinbar ist; für sich allein definiert
sie aber keine Zeitentwicklung, keinen Rückkehrmechanismus und keine
Spektrallücke der Navier--Stokes-Dynamik. Der Ein-Parameter-Freiheitsgrad in der Kolmogorov--Heisenberg-Schließung ($\Phi_j \equiv \alpha$)
entspricht einer einzelnen flachen Richtung in der $\Pi$-Faser des
erweiterten Raums $(\mathbf{E}, \Pi)$, die die strikte
Positivität von $D^2_{\mathbf{E}}\mathcal{F}$ nicht beeinträchtigt. Diese Struktur ist eine
Instanz des in~\cite{FST-Unified2026} diskutierten Resolventen-Dominanz-Musters zweiter Ordnung.
\end{proof}

\begin{remark}[Diagonale Hesse-Schranke]\label{rem:hessian}
Die Hesse-Matrix von $\mathcal{F}$ am K41-Minimierer ist diagonal mit Einträgen
\[
  \frac{\partial^2\mathcal{F}}{\partial E_j^2}
  \;=\; \frac{1}{E_j^*}.
\]
Damit gilt für den kleinsten Eigenwert
\[
  \min_j\frac{1}{E_j^*}
  \;=\;
  \frac{1}{C_K\,\varepsilon^{2/3}\,k_1^{-5/3}\,\Delta k_1}
  \;=\;
  \frac{1}{C_K\,\varepsilon^{2/3}(\ln 2)\,k_1^{-2/3}}
  \;>\; 0.
\]
Dies liefert die lokale gewichtete Stabilitätsskala des Variationsfunktionals.
Allein daraus folgt jedoch kein Spektrallückensatz für die
Navier--Stokes-Dynamik.
\end{remark}

\begin{remark}[Zur Rolle des Referenzspektrums]\label{rem:circularity}
Ein natürlicher Einwand ist, dass das KL-Funktional $\mathcal{F}$
\emph{konstruktionsbedingt} bei $\mathbf{E}^*$ minimiert wird, ungeachtet
des physikalischen Gehalts von~$\mathbf{E}^*$. In der Tat ist für jede feste
Schließung $\Pi$ das unrestringierte Minimum von $\mathcal{F}$ stets
$\mathbf{E}^*$. Der nichttriviale Inhalt von Theorem~\ref{thm:K41} ist
daher nicht, dass $\mathcal{F}$ sein Minimum bei $\mathbf{E}^*$ erreicht
--- das ist eine Konsequenz der KL-Konstruktion --- sondern vielmehr, dass:
\begin{enumerate}
\item Unter allen stationären Profilen $\mathbf{E}^{(\Pi)}$ (eines pro Schließung
  $\Pi \in \mathscr{A}$) einzig das K41-Profil $\mathbf{E}^*$
  die Bedingung konstanten Flusses \emph{simultan erfüllt} und
  $\mathcal{F} = 0$ erreicht. Jedes andere schließungskompatible Profil
  weist $\mathcal{F} > 0$ auf.
\item Die Referenz $\mathbf{E}^*$ innerhalb der normalisierten
  Kolmogorov--Heisenberg-Schließung nicht beliebig ist: Die Relation
  $\Pi_j=\alpha k_j E_j^{3/2}=\varepsilon$ fixiert die Skalierung
  $E_j\propto k_j^{-2/3}$, also für die Spektraldichte
  $E(k)\propto k^{-5/3}$. Breitere oder anders normalisierte
  Schließungsfamilien können natürlich um andere Referenzen konstruiert
  werden; das Theorem ist daher eine Selektionsaussage innerhalb der
  angegebenen K41-normalisierten zulässigen Klasse, kein Beweis, dass keine
  andere Modellkonvention formulierbar ist.
\end{enumerate}
Zusammenfassend: Die physikalische Eingabe ist die Bedingung konstanten Flusses; das KL-Funktional
liefert den Selektionsmechanismus innerhalb von deren Lösungsmenge.
\end{remark}

\begin{remark}[Selbstkonsistenz der K41-Referenz]
\label{rem:k41-selfconsistency}
Das K41-Spektrum wird als Referenzmaß des KL-Funktionals verwendet; sein
Exponent ist aber zugleich selbstkonsistent mit der normalisierten
Konstantfluss-Schließung. Innerhalb des reduzierten lokalen Schalenansatzes
$\Pi_j\sim k_j E_j^{3/2}$ ergibt konstanter Fluss
$E_j\propto \varepsilon^{2/3}k_j^{-2/3}$ und damit
$E(k)\propto \varepsilon^{2/3}k^{-5/3}$. Dies ist ein Konsistenzcheck der
variationellen Referenz, keine unabhängige Herleitung aller
Turbulenzstatistiken.
\end{remark}

\begin{remark}[Heuristische Reaktionsketten-Analogie]\label{rem:nash}
Die Variationsstruktur von Theorem~\ref{thm:K41} erlaubt eine nützliche
Analogie zur Sprache von Reaktionsketten und Spieltheorie. Betrachtet man
den Trägheitsbereich als eine hierarchische \emph{Reaktionskette}: große
Wirbel auf der Skala $k_j$ transferieren Energie an Wirbel auf der Skala
$k_{j+1}$, wobei jede Skala als eine „Art“ agiert, deren Fitness durch ihren
Energieanteil gesteuert wird. Ein mögliches Replikator-Modell auf dem
Simplex der normierten Schalenenergien
$p_j = E_j / \sum_\ell E_\ell$ nimmt die Form
\[
  \dot{p}_j
  \;=\;
  p_j\,\bigl(g_j(\mathbf{p}) - \bar{g}(\mathbf{p})\bigr)
\]
an, wobei $g_j$ die Netto-Energiegewinnrate auf der Skala $j$ und $\bar{g}$ der
Populationsdurchschnitt ist. Mit einer separat spezifizierten Payoff-Struktur,
die mit dem KL-Funktional kompatibel ist, lässt sich das K41-Profil als
Nash-artiger stationärer Punkt darstellen, und $\mathcal{F}$ kann als
Kandidat für eine Ljapunow-Funktion dienen. Dies ist nicht Teil von
Theorem~\ref{thm:K41}: Hier wird keine Replikator-Dynamik, kein Satz zur
asymptotischen Stabilität und kein Rückkehrmechanismus der
Navier--Stokes-Dynamik bewiesen.
\end{remark}

% ===================================================================
\section{Robustheits-Leitplanken und formale Penalty-Struktur}\label{sec:robustness}

Die strenge Flussbedingung von Theorem~\ref{thm:K41} erfordert
punktweise $\Pi_j = \varepsilon$ über den Trägheitsbereich hinweg.
In der Praxis zeigen DNS-Daten kleinskalige
Fluktuationen um diese Idealisierung. Wir halten zwei Geltungsbereichsgrenzen
fest: Relaxierung allein ist trivial, solange $E^*$ zulässig bleibt, und die
Minimax-Notation ist nur eine formale Penalty-Darstellung.

\begin{proposition}[Relaxierte Flussbedingung: Geltungsbereich]
\label{prop:dfc-slack}
Sei $\mathcal{A}_{\textup{slack}}$ eine beliebige relaxierte
flusszulässige Menge, welche das exakte K41-Profil $\mathbf{E}^*$ enthält.
Wenn die Zielfunktion dasselbe KL-Funktional $\mathcal{F}$ bleibt, dann ist
der Minimierer über $\mathcal{A}_{\textup{slack}}$ weiterhin exakt
$\mathbf{E}^*$:
\begin{equation}\label{eq:dfc-slack-bound}
  \arg\min_{\mathbf{E}\in\mathcal{A}_{\textup{slack}}}
  \mathcal{F}[\mathbf{E}]
  = \mathbf{E}^*.
\end{equation}
Ein nichttriviales Robustheitstheorem für endliche Reynolds-Zahlen entsteht
daher nicht allein dadurch, dass die Flussungleichungen relaxiert werden,
während dieselbe referenzzentrierte Zielfunktion erhalten bleibt. Dafür
bräuchte man zusätzlich einen Datenanpassungsterm, eine gestörte Referenz
oder empirische Schranken, die das exakte Profil aus der Vergleichsklasse
ausschließen.
\end{proposition}

\begin{proof}
Für alle positiven Schalenspektren gilt $\mathcal{F}[\mathbf{E}]\ge0$, und
Gleichheit tritt nur bei $\mathbf{E}=\mathbf{E}^*$ ein. Liegt
$\mathbf{E}^*$ in $\mathcal{A}_{\textup{slack}}$, dann ist das relaxierte
Infimum höchstens $\mathcal{F}[\mathbf{E}^*]=0$ und mindestens $0$.
Jeder relaxierte Minimierer mit Zielfunktionswert null ist daher
$\mathbf{E}^*$. Dies beweist die Behauptung und erklärt, warum die ältere
Formulierung über ``Slack-Minimierer'' zu stark war: Das exakte
Referenzprofil bleibt zulässig.
\end{proof}

\begin{remark}
Proposition~\ref{prop:dfc-slack} ist eine Leitplanke, kein positives
Robustheitstheorem: Robustheit bei endlicher Reynolds-Zahl muss mit einer
gestörten Referenz, einem empirischen Datenanpassungsterm oder einer
Vergleichsklasse formuliert werden, die das exakte K41-Profil nicht mehr
enthält. Ohne eine solche zusätzliche Zutat bleibt das KL-Minimum
konstruktionsbedingt exakt bei der Referenz.
\end{remark}

\begin{proposition}[Formale Penalty-Darstellung]
\label{prop:minimax}
In einer endlichen Schalentrunkierung mit
$E_{\mathrm{tot}}=\sum_j E_j^*$ besitzt das entkoppelte Penalty-Problem
\begin{equation}\label{eq:minimax}
  \mathbf{E}^*
  \;=\;
  \arg\min_{\mathbf{E} \in \mathcal{E}}\;
  \max_{\boldsymbol{\Pi} \in \mathcal{A}}\;
  \Bigl[
    \mathcal{F}[\mathbf{E}]
    \;-\;
    \lambda \sum_j (\Pi_j - \varepsilon)^2
  \Bigr],
\end{equation}
den Sattelpunkt
$(\mathbf{E}^*,(\varepsilon,\ldots,\varepsilon))$, wobei $\mathcal{E}$
die Menge der nichtnegativen Energiespektren mit fester Gesamtenergie und
$\mathcal{A}$ eine beschränkte konvexe Menge zulässiger Flussprofile ist.
Da diese Auszahlung $\mathbf{E}$ und $\boldsymbol{\Pi}$ nicht koppelt, ist
sie eine formale Penalty-Darstellung der harten Nebenbedingung, kein
dynamischer Zwei-Spieler-Satz für den Navier--Stokes-Transfer.
\end{proposition}

\begin{proof}
Wir überprüfen die endlichdimensionale Minimax-Struktur, identifizieren den
Sattelpunkt und machen explizit, wo die Darstellung entkoppelt ist.

\smallskip\noindent
\textbf{Schritt 1: Strategieräume.}\;
Definiere die Lagrange-bestrafte Auszahlung
\[
  L(\mathbf{E}, \boldsymbol{\Pi})
  \;:=\;
  \mathcal{F}[\mathbf{E}]
  \;-\;
  \lambda \sum_{j \in \mathcal{J}} (\Pi_j - \varepsilon)^2.
\]
Die Spektrumsvariable liegt in
$\mathcal{E} := \{\mathbf{E} \in \mathbb{R}^N_{>0} :
\sum_j E_j = E_{\mathrm{tot}}\}$, dem Simplex der positiven
Energiespektren mit vorgegebener Gesamtenergie $E_{\mathrm{tot}} > 0$ und
$N = |\mathcal{J}|$ Schalen. Die Flussvariable ist
$\boldsymbol{\Pi} = (\Pi_j)_{j=1}^{N-1}$ und liegt in
$\mathcal{A} := \{\boldsymbol{\Pi} \in \mathbb{R}^{N-1} :
0 \le \Pi_j \le \Pi_{\max}\}$ für eine feste obere Schranke
$\Pi_{\max} > \varepsilon$ (physikalisch ist der Fluss im Trägheitsbereich
durch $O(\varepsilon)$ nach oben beschränkt).

Die Menge $\mathcal{E}$ ist eine kompakte konvexe Teilmenge von $\mathbb{R}^N$
(sie ist der Schnitt des offenen positiven Orthanten mit der
Hyperebene $\sum_j E_j = E_{\mathrm{tot}}$; die Kompaktheit gilt, da
$E_j \ge 0$ und $E_j \le E_{\mathrm{tot}}$ für jedes $j$, und wir betrachten
den Abschluss in $\mathbb{R}^N_{\ge 0}$, wobei $L$ stetig auf den Rand
fortsetzbar ist). Die Menge $\mathcal{A}$ ist eine kompakte
konvexe Teilmenge von $\mathbb{R}^{N-1}$ (ein abgeschlossener Quader).

\smallskip\noindent
\textbf{Schritt 2: Konvexität--Konkavität.}\;
Wir verifizieren:
\begin{enumerate}[label=(\alph*)]
  \item \emph{$L$ ist konvex in $\mathbf{E}$ für jedes feste
    $\boldsymbol{\Pi}$.}\;
    Der Term $\mathcal{F}[\mathbf{E}] = \sum_j [E_j\ln(E_j/E_j^*)
    - E_j + E_j^*]$ ist strikt konvex in $\mathbf{E}$, da seine
    Hesse-Matrix $\mathrm{diag}(1/E_j)$ für
    $\mathbf{E} > 0$ positiv definit ist. Der Strafterm
    $-\lambda\sum_j(\Pi_j - \varepsilon)^2$ ist unabhängig
    von $\mathbf{E}$. Daher ist $L(\cdot, \boldsymbol{\Pi})$
    strikt konvex.
  \item \emph{$L$ ist konkav in $\boldsymbol{\Pi}$ für jedes feste
    $\mathbf{E}$.}\;
    Für festes $\mathbf{E}$ ist $\mathcal{F}[\mathbf{E}]$ konstant
    bzgl. $\boldsymbol{\Pi}$. Die Abbildung $\boldsymbol{\Pi} \mapsto
    -\lambda\sum_j(\Pi_j - \varepsilon)^2$ ist strikt konkav
    (ihre Hesse-Matrix ist $-2\lambda\,I_{N-1}$, negativ definit).
    Daher ist $L(\mathbf{E}, \cdot)$ strikt konkav.
\end{enumerate}

\smallskip\noindent
\textbf{Schritt 3: Anwendung des Sion'schen Minimax-Theorems.}\;
Nach dem Minimax-Theorem von Sion (M.~Sion, \emph{Pacific J.\ Math.}\ 8,
1958, 171--176) gilt: Ist $X$ eine kompakte konvexe Teilmenge eines topologischen
Vektorraums, $Y$ eine konvexe Teilmenge eines topologischen Vektorraums,
und $f: X \times Y \to \mathbb{R}$ derart, dass $f(\cdot, y)$ für
jedes $y$ konvex und unterhalbstetig ist, und $f(x, \cdot)$ für
jedes $x$ konkav und oberhalbstetig ist, so gilt
$\min_X \sup_Y f = \sup_Y \min_X f$.

Hier ist $X = \mathcal{E}$ (kompakt konvex), $Y = \mathcal{A}$ (kompakt
konvex, also insbesondere konvex), und $L$ ist stetig (somit
sowohl unter- als auch oberhalbstetig) und erfüllt die
Konvexitäts-Konkavitäts-Bedingungen aus Schritt~2. Kompaktheit beider
Mengen und Stetigkeit stellen sicher, dass $\min$ und $\max$ angenommen werden.
Daher gilt:
\begin{equation}\label{eq:sion-applied}
  \min_{\mathbf{E} \in \mathcal{E}}\;
  \max_{\boldsymbol{\Pi} \in \mathcal{A}}\;
  L(\mathbf{E}, \boldsymbol{\Pi})
  \;=\;
  \max_{\boldsymbol{\Pi} \in \mathcal{A}}\;
  \min_{\mathbf{E} \in \mathcal{E}}\;
  L(\mathbf{E}, \boldsymbol{\Pi}),
\end{equation}
und ein Sattelpunkt $(\mathbf{E}^\dagger, \boldsymbol{\Pi}^\dagger)$
existiert, welcher
$L(\mathbf{E}^\dagger, \boldsymbol{\Pi})
\le L(\mathbf{E}^\dagger, \boldsymbol{\Pi}^\dagger)
\le L(\mathbf{E}, \boldsymbol{\Pi}^\dagger)$
für alle $\mathbf{E} \in \mathcal{E}$,
$\boldsymbol{\Pi} \in \mathcal{A}$ erfüllt.

\smallskip\noindent
\textbf{Schritt 4: Identifikation des Sattelpunkts.}\;
Am Sattelpunkt müssen die Bedingungen erster Ordnung erfüllt sein.

\emph{Stationarität in $\boldsymbol{\Pi}$:}\;
Für festes $\mathbf{E}^\dagger$ ist die Abbildung
$\boldsymbol{\Pi} \mapsto L(\mathbf{E}^\dagger, \boldsymbol{\Pi})$
strikt konkav und differenzierbar. Der innere Maximierer
erfüllt
\[
  \frac{\partial L}{\partial \Pi_j}
  \;=\;
  -2\lambda\,(\Pi_j - \varepsilon)
  \;=\; 0
  \qquad\Longrightarrow\qquad
  \Pi_j^\dagger = \varepsilon \quad\text{für alle } j.
\]
(Diese innere Lösung ist zulässig, da
$0 < \varepsilon < \Pi_{\max}$.)

\emph{Stationarität in $\mathbf{E}$:}\;
Für festes $\boldsymbol{\Pi}^\dagger = (\varepsilon, \dots, \varepsilon)$
verschwindet der Strafterm und
$L(\mathbf{E}, \boldsymbol{\Pi}^\dagger) = \mathcal{F}[\mathbf{E}]$.
Unter der Nebenbedingung $\mathbf{E} \in \mathcal{E}$ (die Energiebedingung
$\sum_j E_j = E_{\mathrm{tot}}$) minimieren wir $\mathcal{F}$ über die
Methode der Lagrange-Multiplikatoren. Die Stationaritätsbedingung lautet
\[
  \frac{\partial \mathcal{F}}{\partial E_j}
  \;=\;
  \ln\frac{E_j}{E_j^*}
  \;=\; \mu
  \qquad\text{für alle } j \in \mathcal{J},
\]
wobei $\mu$ der Multiplikator für die Gesamtenergiebedingung ist.
Dies ergibt $E_j = e^\mu\,E_j^*$ für alle $j$, d.h.\ der Minimierer ist
proportional zum K41-Referenzspektrum. Die Konstante $e^\mu$ wird
durch die Gesamtenergiebedingung
$\sum_j E_j = e^\mu \sum_j E_j^* = E_{\mathrm{tot}}$ festgelegt. Nimmt man
$E_{\mathrm{tot}} = \sum_j E_j^*$ (welche die Gesamtenergie des
K41-Spektrums ist), erhält man $\mu = 0$ und damit
$E_j^\dagger = E_j^*$ für alle $j$.

Daher ist der Sattelpunkt
$(\mathbf{E}^\dagger, \boldsymbol{\Pi}^\dagger) =
(\mathbf{E}^*, (\varepsilon, \dots, \varepsilon))$:
das K41-Spektrum gepaart mit dem konstanten Flussprofil.

\smallskip\noindent
\textbf{Schritt 5: Eindeutigkeit.}\;
Der Sattelpunkt ist eindeutig, weil $L$ \emph{strikt} konvex in
$\mathbf{E}$ und \emph{strikt} konkav in $\boldsymbol{\Pi}$ ist.
Tatsächlich, wenn $(\tilde{\mathbf{E}}, \tilde{\boldsymbol{\Pi}})$ ein
anderer Sattelpunkt wäre, würde die Sattelpunkt-Ungleichung
$L(\tilde{\mathbf{E}}, \boldsymbol{\Pi}^\dagger) \ge
L(\mathbf{E}^\dagger, \boldsymbol{\Pi}^\dagger)$ und
$L(\mathbf{E}^\dagger, \boldsymbol{\Pi}^\dagger) \le
L(\tilde{\mathbf{E}}, \boldsymbol{\Pi}^\dagger)$ liefern, also
$L(\tilde{\mathbf{E}}, \boldsymbol{\Pi}^\dagger) =
L(\mathbf{E}^\dagger, \boldsymbol{\Pi}^\dagger)$. Aufgrund strikter
Konvexität in $\mathbf{E}$ erzwingt dies
$\tilde{\mathbf{E}} = \mathbf{E}^\dagger = \mathbf{E}^*$.
Ebenso erzwingt strikte Konkavität
$\tilde{\boldsymbol{\Pi}} = \boldsymbol{\Pi}^\dagger$.

\smallskip\noindent
\textbf{Schritt 6: Geltungsbereich der Penalty-Darstellung.}\;
Die Penalty-Formulierung~\eqref{eq:minimax} ist ein entkoppeltes
Buchhaltungsinstrument für die Konstantflussbedingung
$\Pi_j=\varepsilon$. Für feste endlichdimensionale Schalenvektoren zwingt
der Term $-\lambda(\Pi_j-\varepsilon)^2$ jede maximierende Flussvariable
im Grenzfall $\lambda\to\infty$ gegen $\varepsilon$; in diesem Grenzfall
hat der formale Sattelpunkt denselben KL-Minimierer wie das
konstantflussbeschränkte Problem aus Theorem~\ref{thm:K41}.

Dies ist keine Äquivalenz zu einem closure-gekoppelten harten Problem
$\Pi_j[\mathbf{E}]=\varepsilon$ für beliebige feste Schließungen. Eine
solche Aussage würde Kompaktheit, Constraint-Qualification und
Existenzsätze für die gewählte Schließungsfamilie erfordern. Die
vorliegende Minimax-Aussage ist nur das endlichdimensionale
Konstantfluss-Penalty-Buchhaltungslemma, das dieses Paper verwendet.
\end{proof}

\begin{remark}
Die Minimax-Notation ist als Buchhaltung für die Konstantfluss-Penalty
nützlich, zeigt aber für sich genommen keine nichttriviale
spieltheoretische Struktur der turbulenten Kaskade. Ein echtes
Zwei-Spieler-Modell müsste eine gekoppelte Auszahlung besitzen, in der das
gewählte Flussprofil die zulässige Energieentwicklung oder das
Dissipationsfunktional verändert. Die vorliegende Formulierung hält nur
fest, dass KL-Minimierer und Konstantfluss-Penalty denselben formalen
Sattelpunkt besitzen.
\end{remark}

\begin{remark}[Lokale gewichtete $\ell^2$-Stabilität aus der Hesse-Schranke]
\label{rem:talagrand}
Das Funktional $\mathcal{F}[\mathbf{E}]$ ist lokal gleichmäßig konvex um
den K41-Minimierer $\mathbf{E}^*$ mit diagonaler Hesse-Matrix
$D^2\mathcal{F}|_{\mathbf{E}^*} = \mathrm{diag}(1/E_j^*)$. Bleiben die
Verhältnisse $E_j/E_j^*$ in einem festen kompakten Intervall
$[R^{-1},R]$, liefert Taylor-Entwicklung eine vom Verhältnis abhängige
gewichtete Stabilitätsschranke
\[
  c_R\sum_j \frac{(E_j - E_j^*)^2}{E_j^*}
  \;\le\; \mathcal{F}[\mathbf{E}]
\]
für ein $c_R>0$. Dies ist die hier verwendete Stabilitätsaussage. Die
stärkere globale Ungleichung
$x\ln(x/y)-x+y\ge (x-y)^2/(2y)$ ist für große $x/y$ falsch und wird nicht
verwendet.
\end{remark}

\begin{remark}[Lokales UCU3 in der abstrakten Drei-Lemma-Architektur]
\label{rem:local-ucu3}
Die gewichtete $\ell^2$-Schranke in Bemerkung~\ref{rem:talagrand} kann als
die K41-Instanziierung des abstrakten \emph{Quadratischen-Wachstums-Axioms}
\textnormal{UCU3} der Drei-Lemma-Architektur für strikt konvexe
Variationsfunktionale gelesen werden (Universal Convexity Uniqueness; vgl.\ die
begleitenden Architektur-Notizen \path{CORE/UCU_FORMAL_DEFINITIONS.md}~\S5
und \path{CORE/ABSTRACT_THREE_LEMMA_UCU.md}~\S4.2.3). Definiert man die
gewichtete Norm $\|h\|^2_*\,:=\,\sum_j h_j^2/E_j^*$, lautet die Schranke
\[
  \mathcal{F}[\mathbf{E}] - \mathcal{F}[\mathbf{E}^*]
  \;\ge\;
  c_R\,\|\mathbf{E}-\mathbf{E}^*\|_*^{\,2},
\]
d.h.\ \emph{lokales} UCU3 gilt mit einer verhältnisabhängigen Konstante
$c_R>0$ in der gewichteten
$\ell^2(1/E_j^*)$-Norm. Zusammen mit der strikten Konvexität (UCU1, unser
Theorem~\ref{thm:K41} Teil~(ii)) und der Existenz des K41-Minimierers
(UCU2, Theorem~\ref{thm:K41} Teil~(i)) schließt dies die lokale
Drei-Lemma-Architektur für K41 im Standardsinne der konvexen
Analysis.

\smallskip\noindent\textit{Was offen bleibt (globales UCU3, ehrliche Aussage des Geltungsbereichs).} Die obige Schranke ist eine lokale Aussage um $\mathbf{E}^*$
(Taylor-Entwicklung bis zur zweiten Ordnung, kontrollierter Restterm über
ein festes kompaktes Verhältnisintervall). Die stärkere globale Ungleichung
$x\ln(x/y)-x+y\ge (x-y)^2/(2y)$ ist für große $x/y$ falsch und wird hier
nicht verwendet. Ein \emph{globales} UCU3 im
Sinne einer Transport-Informations-Ungleichung vom Talagrand-Typ
$\mathrm{KL}(\mu\|\mu_{K41}) \ge (\rho_{K41}/2)\,W_2^2$ für ein Wahrscheinlichkeitsmaß $\mu_{K41}$ auf Kaskaden-Konfigurationen wird, über die
Otto--Villani-Route~\cite{OttoVillani2000}, auf eine log-Sobolev-Ungleichung
für $\mu_{K41}$ mit positiver Konstante $\rho_{K41}$ zurückgeführt. Die Existenz und
explizite Form einer solchen Konstanten für K41-Kaskaden --- insbesondere ihre
Asymptotik hinsichtlich Reynolds-Zahl und Kaskadentiefe --- ist unseres Wissens nach
in der Standard-Turbulenzliteratur nicht dokumentiert und bleibt ein ungelöstes
Problem in der mathematischen statistischen Mechanik. Eine strukturierte Reduktions-Notiz
findet sich in \texttt{TALAGRAND\_CONSTANT\_K41.md} (begleitende
Arbeitsnotiz).
\end{remark}

% ===================================================================
\section{Intermittenzkorrekturen}\label{sec:intermittency}

Die K41-Theorie sagt voraus, dass die longitudinale Strukturfunktion $p$-ter Ordnung wie folgt skaliert:
\begin{equation}\label{eq:Sp_K41}
  S_p(\ell) \;:=\;
  \bigl\langle |\delta_\ell u|^p \bigr\rangle
  \;\sim\; (\varepsilon\,\ell)^{p/3},
  \qquad
  \zeta_p^{\mathrm{K41}} = \frac{p}{3}.
\end{equation}
Experimente zeigen systematische Abweichungen auf: $\zeta_p < p/3$ für $p > 3$,
eine Signatur der \emph{Intermittenz}.

\subsection{Fluktuationen um das Variationsminimum}

In diesem Variationsrahmen kann Intermittenz modelliert werden, indem man
Fluktuationen der Schalenenergien $\mathbf{E}$ um den K41-Minimierer
$\mathbf{E}^*$ ergänzt. Die Entwicklung von $\mathcal{F}$ bis zur zweiten
Ordnung liefert:
\begin{equation}\label{eq:F_expand}
  \mathcal{F}[\mathbf{E}^* + \delta\mathbf{E}]
  \;=\;
  \mathcal{F}[\mathbf{E}^*]
  \;+\;
  \frac{1}{2}\sum_{j}
    \frac{(\delta E_j)^2}{E_j^*}
  \;+\; O(|\delta\mathbf{E}|^3).
\end{equation}
Die Hesse-Matrix $H_{jj} = 1/E_j^*$ impliziert auf der Ebene der statischen
quadratischen Näherung, dass Fluktuationen auf kleinen Skalen (große $j$,
kleines $E_j^*$) stärker bestraft werden als Fluktuationen auf großen
Skalen. Diese Asymmetrie liefert eine variationelle Modellierungsroute für
Intermittenz; sie ist keine Herleitung der Intermittenz aus der
Navier--Stokes-Dynamik.

\subsection{Verbindung zum K62-Lognormal-Modell}

Wenn wir die Fluktuationen $\delta E_j$ als zusätzlichen Modellansatz als
gaußverteilt in der Variable
$\ln(E_j/E_j^*)$ modellieren, und ferner annehmen, dass Phasen-Korrelationen zwischen
Littlewood--Paley-Blöcken vernachlässigbar sind (sodass sich Schalenenergie-Fluktuationen
direkt in die Skalierung der Strukturfunktionen übersetzen), so ergeben sich die resultierenden Strukturfunktion-Exponenten als
\begin{equation}\label{eq:zeta_p}
  \zeta_p
  \;=\;
  \frac{p}{3}
  \;-\;
  \frac{\mu}{18}\,p(p-3),
\end{equation}
was die verfeinerte Ähnlichkeitshypothese (refined similarity hypothesis)
von K62 (Kolmogorov--Obukhov) mit Intermittenzparameter $\mu$
reproduziert~\cite{K62,Frisch95}. Im formalen Gibbs-Modell
$\exp(-\mathcal{F}/\Theta)$ parametrisiert die effektive Temperatur
$\Theta > 0$ die Stärke der Fluktuationen um das K41-Gleichgewicht. Für
$x_j=\ln(E_j/E_j^*)$ gilt lokal
\[
  \mathcal{F}[\mathbf{E}]
  \;=\;
  \frac{1}{2}\sum_j E_j^* x_j^2 + O(|x|^3),
\]
sodass die quadratische Gibbs-Näherung Log-Energie-Varianzen der Ordnung
$\Theta/E_j^*$ liefert. Der Grenzfall $\Theta\to0$ unterdrückt in diesem
Modell also Intermittenz. Eine Identifikation dieser schalenweisen
Varianzen mit einem universellen K62-Parameter $\mu$ verlangt jedoch eine
zusätzliche Brücke über Kaskadenmaß, refined similarity und
Strukturfunktionen; die Hesse-Skala allein bestimmt kein
skalenunabhängiges $\mu$.

\begin{remark}
Das She--Lévêque-Modell~\cite{SL94} schlägt folgende Alternative vor:
\[
  \zeta_p = \frac{p}{9} + 2\Bigl(1 - \bigl(\tfrac{2}{3}\bigr)^{p/3}\Bigr),
\]
welches experimentelle Daten für große $p$ genauer approximiert. Eine Ableitung hiervon
aus dem Freie-Energie-Rahmen würde erfordern, die Fluktuationen als
Log-Poisson- statt als Log-Normal-Prozess zu modellieren. Dies bleibt ein offenes Problem
innerhalb unseres Ansatzes.
\end{remark}

% ===================================================================
\subsection{Strukturelle Klassifikation}\label{sec:resolvent-cascade}

Die gemeinsame Minimierung aus Theorem~\ref{thm:K41} weist eine dreistufige
Hierarchie auf: (i)~das Energiebudget in führender Ordnung ist neutral (jedes Potenzgesetz-Spektrum ist kinematisch zulässig); (ii)~die Bedingung konstanten Flusses erster Ordnung schränkt das Profil ein, legt es aber nicht
eindeutig fest; (iii)~die strikte Konvexität zweiter Ordnung
von~$\mathcal{F}$ wählt K41 als eindeutigen Minimierer. Dieses Muster --- in~\cite{FST-Unified2026} als \emph{Resolventen-Dominanz zweiter Ordnung} bezeichnet --- hat Analoga in anderen variationellen Selektionsproblemen und wird dort ausführlich diskutiert.

% ===================================================================
\section{Numerische Illustration und DNS-Diagnostik}\label{sec:numerics}

\begin{remark}[Numerische Illustration der KL-Landschaft]\label{rem:numerical-F}
Eine numerische Auswertung von $\mathcal{F}[\mathbf{E}]$ über sechs synthetische Energiespektren illustriert die referenzzentrierte KL-Landschaft: $\mathcal{F}[\mathbf{E}_{\mathrm{K41}}] = 0$ exakt, während $\mathcal{F}[\mathbf{E}_{\mathrm{K62}}] = 2{,}0 \times 10^{-4}$ (Intermittenzkorrektur), $\mathcal{F}[\mathbf{E}_{k^{-2}}] = 1{,}74$ (steil), $\mathcal{F}[\mathbf{E}_{k^{-4/3}}] = 5{,}94$ (flach) und $\mathcal{F}[\mathbf{E}_{\mathrm{thermal}}] = 7{,}55$ (Gleichverteilung). Ein Perturbationstest ($300$ zufällige Störungen $\delta \mathbf{E}$ bei drei Rauschleveln $\|\delta \mathbf{E}\|/\|\mathbf{E}^*\| \in \{0{,}01,\, 0{,}1,\, 1{,}0\}$) liefert $\mathcal{F}[\mathbf{E}^* + \delta \mathbf{E}] > 0$ für \emph{alle} $300$ Stichproben, was die strikte Konvexität numerisch illustriert. Numerische Skripte: \texttt{compute\_F\_spectrum.py}.
\end{remark}

Die Audit-Ansicht in Tabelle~\ref{tab:numerical-diagnostics} macht die
diagnostische Rolle der synthetischen Spektren explizit. Das
Bottleneck-Profil liegt in $\mathcal{F}$ nahe an K41, scheitert aber am
strikten K41-Steigungsgate; thermische und warme Randrichtungen liegen
außerhalb der K41-Diagnose-Wasserlinie.

\begin{table}[ht]
\centering
\small
\setlength{\tabcolsep}{4pt}
\begin{tabular}{|L{0.22\textwidth}|L{0.13\textwidth}|L{0.23\textwidth}|L{0.30\textwidth}|}
\hline
\textbf{Synthetisches Spektrum} & \textbf{$\mathcal{F}[E]$} & \textbf{Gate-Status} & \textbf{Diagnostische Lesart} \\
\hline
K41-Referenz & $0{,}000000$ & DFC1 ja; schwache und K41-Steigung ja & Exakter Referenzminimierer. \\
\hline
K62, $\mu=0{,}03$ & $0{,}000197$ & DFC1 ja; schwache und K41-Steigung ja & K41-nahe Intermittenzkontrolle. \\
\hline
Bottleneck & $0{,}011850$ & DFC1 ja; schwach ja; K41-Steigung nein & Stresskontrolle: kleine KL-Abweichung, striktes Steigungsgate scheitert. \\
\hline
Steil $k^{-2}$ & $1{,}736741$ & DFC1 ja; Steigungsgates nein & Nicht-K41-Steilkontrolle. \\
\hline
Flach $k^{-4/3}$ & $5{,}941487$ & DFC1 ja; Steigungsgates nein & Nicht-K41-Flachkontrolle. \\
\hline
Thermisch flach & $7{,}552260$ & DFC1 nein; Steigungsgates nein & Nicht-Kaskaden-Randfall. \\
\hline
\end{tabular}
\caption{Synthetische Freie-Energie- und Gate-Diagnostik aus \texttt{compute\_F\_spectrum.py}.}
\label{tab:numerical-diagnostics}
\end{table}

Wir schlagen folgendes reproduzierbares Diagnoseprotokoll für DNS-Tests vor:
\begin{enumerate}
\item Berechne die Littlewood--Paley Schalenenergien $E_j$ aus
  einer statistisch stationären DNS bei $R_\lambda \ge 200$.
\item Werte $\mathcal{F}[\mathbf{E}]$ unter Verwendung des gemessenen
  $\mathbf{E}$ und der K41-Referenz $\mathbf{E}^*$ aus, wobei
  $\varepsilon := \nu\,\langle\|\nabla u\|_{L^2}^2\rangle_t$
  die zeitgemittelte viskose Dissipationsrate aus demselben DNS-Lauf ist.
\item Teste, ob $\mathcal{F}$ mit zunehmendem $R_\lambda$ abnimmt
  und dass das kompensierte Verhältnis $\varphi_j = \ln(E_j/E_j^*)$ über den
  Trägheitsbereich hinweg monoton fallend in $j$ ist.
\end{enumerate}

\smallskip\noindent
\textbf{Implementierungshinweise.}\;
(a)~Standard-DNS-Codes geben das isotrope Energiespektrum $E(k)$
via scharfer sphärischer Schalen-Gruppierung (binning) aus, nicht Littlewood--Paley (LP) gefilterte Energien. Für glatte Spektren ist der Unterschied von der Größenordnung $O(\Delta k_j / k_j) = O(\ln 2)$ pro Schale und kann abgeschätzt werden, indem man scharfes und glattes Binning vergleicht.
(b)~Die Schritte 1--3 sollten auf instantane Spektren bei mehreren
statistisch unabhängigen Zeitpunkten innerhalb des stationären Regimes angewendet werden;
das Ensemblemittel von $\mathcal{F}$ und $\varphi_j$ liefert dann
Konfidenzintervalle.

% ===================================================================
\section{Diskussion}\label{sec:discussion}

Die gemeinsame Minimierung über die K41-normalisierte zulässige Klasse ist nicht bloß die Aussage, dass eine KL-Divergenz an ihrer Referenz minimiert wird: Kolmogorovs ursprüngliche Ableitung beruht auf Dimensionsanalyse, während Theorem~\ref{thm:K41} festhält, dass dasselbe Profil unter der normalisierten Konstantfluss-Schließung erreichbar ist und auf der daraus entstehenden gemeinsamen zulässigen Menge der eindeutige Energieminimierer ist. Dies ist eine strukturelle Konsistenz- und Selektionsaussage, keine vollständig intrinsische Herleitung von K41.

\subsection{Die Rolle der Schließungsfamilie}

Die Umformulierung von Theorem~\ref{thm:K41} via der zulässigen Flussfamilie (Definition~\ref{def:admissible}) schwächt die Abhängigkeit von einem einzelnen spezifischen Schließungsmodell ab. Das K41-Spektrum wird nicht durch eine einzelne algebraische Relation selektiert, sondern als der eindeutige Energieminimierer in der gemeinsamen zulässigen Menge der K41-normalisierten lokalen, dimensionskonsistenten, monotonen Flussklasse.
Nichtsdestotrotz erzwingt Axiom~(A2) weiterhin eine strukturelle Einschränkung --- den $k_j\,E_j^{3/2}$-Skalierungsvorfaktor ---, der letztendlich das dimensionale Argument von Kolmogorov einkodiert. Die Ableitung dieses Vorfaktors aus der Kármán--Howarth--Monin-Relation oder aus strengen Schranken an die trilineare Littlewood--Paley-Form bleibt ein wichtiges offenes Problem.

\subsection{Intrinsische Charakterisierung von \texorpdfstring{$E^*$}{E*}}

\begin{remark}[Kármán--Howarth--Monin Charakterisierung]
\label{rem:khm-intrinsic}
Die Kármán--Howarth--Monin (KHM)-Relation
$\partial_t \langle |\delta u|^2 \rangle + \nabla_r \cdot
\langle |\delta u|^2 \delta u \rangle = -4\varepsilon/3
+ 2\nu \Delta_r \langle |\delta u|^2 \rangle$ liefert einen intrinsischen
Konstantfluss-Anker für den Trägheitsbereich. Im stationären
Trägheitsbereich ($\eta \ll r \ll L$) führt sie zum Vier-Fünftel-Gesetz für
die Strukturfunktion dritter Ordnung. Das stützt die Flusseingabe des
Variationsschemas, bestimmt aber nicht allein durch Fourier-Transformation
eindeutig das Energiespektrum zweiter Ordnung. Die Rückgewinnung von
$E(k)\sim \varepsilon^{2/3}k^{-5/3}$ benötigt weiterhin die übliche
Selbstähnlichkeits-, Lokalitäts- oder Schließungsannahme. Die KHM-Relation
ist daher ein Konsistenzcheck der Referenz, kein unabhängiger
Eindeutigkeitsbeweis.
\end{remark}

\subsection{Regimeauswahl und laminar-turbulenter Übergang}

Das Variationsprinzip in Theorem~\ref{thm:K41} wählt das K41-Spektrum \emph{innerhalb} des turbulenten Regimes aus, erklärt jedoch den laminar-turbulenten Übergang selbst nicht. Eine natürliche Erweiterung würde die strikte Konvexität von~$\mathcal{F}$ als Stabilitätskriterium nutzen: Unter allen kinematisch zulässigen Strömungsregimen ist das in der Natur realisierte jenes, dessen spektrale Energieverteilung ein stabiler kritischer Punkt des Freie-Energie-Funktionals ist. Diese Verbindung quantitativ zu erforschen bleibt zukünftiger Arbeit überlassen.

\subsection{Einschränkungen}

Wir fassen die hauptsächlichen Einschränkungen der vorliegenden Arbeit zusammen:
\begin{enumerate}
\item \textbf{Statische Auswahl.}\; Theorem~\ref{thm:K41} etabliert K41 als eindeutigen variationellen Minimierer, adressiert jedoch nicht die dynamische Frage, ob die Navier--Stokes-Entwicklung das Spektrum auf diesen Minimierer zutreibt. Die dynamischen Aspekte --- einschließlich $\mathcal{F}$-Monotonie entlang Trajektorien und anomaler Dissipation --- erfordern zusätzliche Hypothesen zur nichtlinearen Kaskade und werden im Begleitpaper~\cite{Geiger2026b} behandelt.
\item \textbf{Trägheitsbereich-Idealisierung.}\; Das Funktional $\mathcal{F}$ ist über den Trägheitsbereich $\mathcal{J}$ definiert und ignoriert die Antriebs- und Dissipationsbereiche. Randeffekte bei $j_{\mathrm{in}}$ und $j_{\mathrm{dis}}$ (insbesondere der Bottleneck) werden in Konstanten absorbiert, aber nicht quantitativ kontrolliert.
\item \textbf{Intermittenz.}\; Die K62-Verbindung (Abschnitt~\ref{sec:intermittency}) ist qualitativ. Eine rigorose Herleitung der Intermittenzexponenten aus den Hesse-Fluktuationen bleibt offen.
\item \textbf{Eindeutigkeit der Referenz.}\; Während Bemerkung~\ref{rem:circularity} $\mathbf{E}^*$ als natürliche Referenz der K41-normalisierten Klasse identifiziert, würde eine vollkommen intrinsische Charakterisierung von $\mathbf{E}^*$ (ohne \emph{a priori} Rückgriff auf K41) die Argumentation stärken. Bemerkung~\ref{rem:khm-intrinsic} stellt einen Schritt in diese Richtung dar.
\end{enumerate}

\subsection{Erweiterungen und Ausblick}

\begin{itemize}
  \item \textbf{Anomale Dissipation.}\;
    Im Begleitpapier~\cite{Geiger2026b} führen wir eine Hierarchie von
    Entropiekompatibilitäts-Hypothesen für die nichtlineare Kaskade ein --- von einer punktweisen Bedingung über eine Variante mit Skalenfenster (scale-windowed) hin zu einer dualen, falsifizierbaren
    \emph{Downhill Flux Condition} (DFC), die an DNS-Daten verifizierbar ist --- und
    beweisen konditional, dass diese zusammen mit den angegebenen uniformen
    Energiezufuhr-Annahmen im Limes verschwindender Viskosität eine untere
    Schranke für die Dissipationsrate implizieren.
  \item \textbf{Kompressible Turbulenz:}
    Die Erweiterung auf kompressible Strömungen erfordert das Ersetzen der
    $L^2$-basierten Schalenenergien durch dichte-gewichtete Größen und das Modifizieren der
    Schließung, um der Kopplung von akustischen Moden Rechnung zu tragen.
  \item \textbf{Zweidimensionale Turbulenz:}
    Die duale Kaskade (inverse Energiekaskade mit $k^{-5/3}$ und direkte
    Enstrophiekaskade mit $k^{-3}$~\cite{Kraichnan67}) stellt einen
    interessanten Testfall dar, der ein Variationsproblem mit zwei Nebenbedingungen erfordert.
\end{itemize}

% ===================================================================
\subsection{Post-Zookeeper-Grenze des Geltungsbereichs}
\label{subsec:zookeeper-k41-scope}

Der Zookeeper-Transfer verdeutlicht, was diese Arbeit beanspruchen sollte und was nicht. Diese Arbeit ist Paper A in der Turbulenz-Aufspaltung: Sie beweist die variationelle Auswahl des K41-Spektrums innerhalb der angegebenen zulässigen Flussfamilie. Sie beweist keine anomale Dissipation für Navier--Stokes, löst nicht das Onsager-Selektionsproblem und behauptet nicht, dass alle schwachen oder Viskositäts-verschwindenden (vanishing-viscosity) Limites K41 auswählen.

Unter dem aktualisierten CORE-Status ist der Transfer lediglich RFEP v1.8
Normalform-Disziplin: Zookeeper schließt Riemann $C^{(2)}$/MS2 im
CCM/Fourier-Zweig, aber daraus folgt keine automatische Turbulenz-Aussage aus
jener Schließung. Das vorliegende Theorem bleibt daher exakt das reine
K41-normalisierte Variationsselektions-Theorem.

Diese Fragen gehören zu einem separaten Paper B. Dort sind die relevanten Objekte
das duale/fensterbasierte DFC-Flussdefizit, das anomale Dissipationsmaß,
gewichtete schalengemittelte Formen des $4/5$-Gesetzes sowie Stresstests aus konvexer Integration und
Vanishing-Viscosity-Konstruktionen. Diese Trennung explizit zu halten, ist eine Stärke: Das statische Variationsresultat in der vorliegenden Arbeit bleibt rein,
während die konditionale Defekttheorie ihre eigene natürliche Fortsetzung hat.

\subsection*{Klassifikation der Resultate}

\begin{center}
\small
\begin{tabular}{|p{0.52\textwidth}|p{0.20\textwidth}|p{0.20\textwidth}|}
\hline
\textbf{Resultat} & \textbf{Typ} & \textbf{Referenz} \\
\hline
\hline
\multicolumn{3}{|l|}{\textit{Theoreme (rigorose statische Variationsaussagen):}} \\
\hline
$\mathcal{F}[\mathbf{E}] \geq 0$, $= 0$ gdw.\ $\mathbf{E} = \mathbf{E}^*$ & Theorem & Thm.~\ref{thm:K41} \\
K41 relativer gemeinsamer Minimierer in der K41-normalisierten Klasse & Theorem & Thm.~\ref{thm:K41} \\
Kanonischer Flussvorfaktor im reduzierten inertialen Schalenansatz & Proposition & Prop.~\ref{prop:A2-derived} \\
\hline
\multicolumn{3}{|l|}{\textit{Geltungsbereich-Leitplanken und formale Umschreibungen:}} \\
\hline
Relaxierter Flussspielraum ist trivial, solange $E^*$ nicht ausgeschlossen oder perturbiert wird & Leitplanke & Prop.~\ref{prop:dfc-slack} \\
Formale Konstantfluss-Penalty-Buchhaltung & Formale Darstellung & Prop.~\ref{prop:minimax} \\
\hline
\multicolumn{3}{|l|}{\textit{Numerisch illustriert:}} \\
\hline
$\mathcal{F}[\mathbf{E}_{\mathrm{K41}}] = 0$ (Minimierer) & Exakt & Rem.~\ref{rem:numerical-F} \\
$\mathcal{F} > 0$ für alle 300 gesampelten Perturbationen & Numerische Illustration & Rem.~\ref{rem:numerical-F} \\
\hline
\multicolumn{3}{|l|}{\textit{DNS-Diagnostiken:}} \\
\hline
$\mathcal{F}[\mathbf{E}_{\mathrm{DNS}}] \to 0$ für $R_\lambda \to \infty$ & Diagnostisches Ziel & Abschn.~\ref{sec:numerics} \\
\hline
\end{tabular}
\end{center}

% ===================================================================
\section*{KI-Offenlegung}
Bei der Erstellung dieses Manuskripts wurden KI-gestützte Tools (Claude von Anthropic und OpenAI Codex) in unterstützender Funktion verwendet, insbesondere für Literaturprüfungen, Textstrukturierung und redaktionelle Verfeinerung. Die konzeptionelle Gestaltung, wissenschaftliche Argumentation und die Verantwortung für den gesamten Inhalt liegen ausschließlich beim Autor.

% ===================================================================
\begin{thebibliography}{99}

\bibitem{BonnansShapiro00}
J.\,F.~Bonnans and A.~Shapiro,
\emph{Perturbation Analysis of Optimization Problems},
Springer Series in Operations Research, Springer, New York, 2000.
\href{https://doi.org/10.1007/978-1-4612-1394-9}{doi:10.1007/978-1-4612-1394-9}.

\bibitem{Frisch95}
U.~Frisch,
\emph{Turbulence: The Legacy of A.\,N.\ Kolmogorov},
Cambridge University Press, 1995.
\href{https://doi.org/10.1017/CBO9781139170666}{doi:10.1017/CBO9781139170666}.

\bibitem{FST-Unified2026}
L.~Geiger, \emph{Functional Stability Theory---Physical Instantiation via the Renormalized Free Energy Principle},
Zenodo, 2026.
\href{https://doi.org/10.5281/zenodo.19036190}{doi:10.5281/zenodo.19036190}.

\bibitem{Geiger2026b}
L.~Geiger,
\emph{Anomalous Dissipation and the Turbulent Cascade: A Variational Free-Energy Characterisation of Kolmogorov Scaling},
Zenodo, 2026.
\href{https://doi.org/10.5281/zenodo.19056813}{doi:10.5281/zenodo.19056813}.

\bibitem{K41a}
A.\,N.~Kolmogorov,
\emph{The local structure of turbulence in incompressible viscous fluid
for very large Reynolds numbers},
Doklady Akad.\ Nauk SSSR \textbf{30} (1941), 299--303.
Reprinted in Proc.\ R.\ Soc.\ Lond.\ A \textbf{434} (1991), 9--13.
\href{https://doi.org/10.1098/rspa.1991.0075}{doi:10.1098/rspa.1991.0075}.

\bibitem{K41b}
A.\,N.~Kolmogorov,
\emph{Dissipation of energy in the locally isotropic turbulence},
Doklady Akad.\ Nauk SSSR \textbf{32} (1941), 16--18.
Reprinted in Proc.\ R.\ Soc.\ Lond.\ A \textbf{434} (1991), 15--17.
\href{https://doi.org/10.1098/rspa.1991.0076}{doi:10.1098/rspa.1991.0076}.

\bibitem{K62}
A.\,N.~Kolmogorov,
\emph{A refinement of previous hypotheses concerning the local structure
of turbulence in a viscous incompressible fluid at high Reynolds number},
J.~Fluid Mech.\ \textbf{13} (1962), 82--85.
\href{https://doi.org/10.1017/S0022112062000518}{doi:10.1017/S0022112062000518}.

\bibitem{Kraichnan67}
R.\,H.~Kraichnan,
\emph{Inertial ranges in two-dimensional turbulence},
Phys.\ Fluids \textbf{10} (1967), 1417--1423.
\href{https://doi.org/10.1063/1.1762301}{doi:10.1063/1.1762301}.

\bibitem{MY75}
A.\,S.~Monin and A.\,M.~Yaglom,
\emph{Statistical Fluid Mechanics}, Vols.~1--2,
MIT Press, 1971--1975.

\bibitem{OttoVillani2000}
F.~Otto and C.~Villani,
\emph{Generalization of an inequality by Talagrand and links with the logarithmic Sobolev inequality},
J.\ Funct.\ Anal.~\textbf{173} (2000), 361--400.
\href{https://doi.org/10.1006/jfan.1999.3557}{doi:10.1006/jfan.1999.3557}.

\bibitem{Pope00}
S.\,B.~Pope,
\emph{Turbulent Flows},
Cambridge University Press, 2000.
\href{https://doi.org/10.1017/CBO9780511840531}{doi:10.1017/CBO9780511840531}.

\bibitem{Rockafellar70}
R.\,T.~Rockafellar,
\emph{Convex Analysis},
Princeton Mathematical Series, vol.~28,
Princeton University Press, 1970.
\href{https://doi.org/10.1515/9781400873173}{doi:10.1515/9781400873173}.

\bibitem{SL94}
Z.-S.~She and E.~L\'ev\^eque,
\emph{Universal scaling laws in fully developed turbulence},
Phys.\ Rev.\ Lett.\ \textbf{72} (1994), 336--339.
\href{https://doi.org/10.1103/PhysRevLett.72.336}{doi:10.1103/PhysRevLett.72.336}.

\end{thebibliography}

\end{document}
